\documentclass[conference]{IEEEtran}
\usepackage[T1]{fontenc}
\usepackage{booktabs}
\usepackage{url}

\begin{document}

\title{How Do People Actually Prompt? Prompt Smells in 493 Real ChatGPT Requests}

\author{\IEEEauthorblockN{Md Habibur Rahman}
\IEEEauthorblockA{MIT Program, Batch 7, Institute of Information Technology\\
Email: shalin.rahman@gmail.com}}

\maketitle

\begin{abstract}
People write prompts every day, but we know little about how well they write them
outside the lab. We took 493 English user prompts from WildChat, a public log of
real ChatGPT conversations, and asked an LLM to label each one with ten
\emph{prompt smells}: signs that a request is hard to answer well. Seven in ten
prompts (71.4\%) had at least one smell. Missing context was by far the most common
problem (58.4\%), followed by an unstated output format (18.7\%) and references
such as ``this'' or ``it'' that point to nothing in the prompt (16.6\%). Smells
were most frequent at the two ends of the length scale. Very short prompts were
mostly vague, while very long prompts tended to be bloated or to set rules that
clash. We release the labeled data and the tool so others can check and extend
the results.
\end{abstract}

\begin{IEEEkeywords}
prompt engineering, prompt smells, large language models, WildChat, empirical study
\end{IEEEkeywords}

\section{Introduction}
Most advice on writing prompts comes from guides and from lab studies with set tasks.
We know much less about how people write prompts when nobody is watching.
WildChat~\cite{wildchat} closes part of that gap. It holds about a million real
conversations that users had with ChatGPT in exchange for free access.

In code, a \emph{smell} is a sign that something may be wrong, even if the program
still runs~\cite{fowler}. We borrow the idea for prompts. A prompt smell is a
pattern that makes a request harder for a model to answer well, such as a missing
goal, an unclear ``it'', or two rules that cannot both hold.

We ask two questions.
\textbf{RQ1:} How common are prompt smells in real English prompts, and which ones
come up most?
\textbf{RQ2:} Do smell rates change with prompt length?

This short paper reports early results from an ongoing study. Our contributions
are (1)~a catalog of ten prompt smells with plain definitions, (2)~smell rates for
493 real prompts, broken down by length, and (3)~an open tool and labeled data set
that anyone can rerun.

\section{Related Work}
White et al.~\cite{white} collect prompt patterns that work, and the Prompt
Report~\cite{promptreport} surveys prompting methods. Both describe what good
prompts look like rather than what users really write. Zamfirescu-Pereira
et al.~\cite{johnny} watched non-experts write prompts and found they often
give up after one failed try and lean on human-style instructions. Our study looks
at the same problem at a larger scale, using logs instead of a lab. We use an
LLM as the labeler, following the LLM-as-a-judge line of work~\cite{judge}.

\section{Method}
\subsection{Data and sampling}
We used one of the six training shards of WildChat (\texttt{train-00003-of-00006},
88,238 conversations). To avoid taking only the earliest conversations in the
file, we kept a conversation when a hash of its id fell in the lowest 1\%. This
gives a fixed, repeatable sample spread across the whole shard. From it we took
every user turn tagged as English that was not empty, 1,003 prompts in all.
The labeler works through them in file order. This paper covers the first 493
prompts it labeled. They come from 238 conversations held between 2 and 15 July
2023, and every one of them was sent to GPT-3.5-Turbo. The median prompt is 139
characters long. Labeling of the remaining prompts is still running.

\subsection{Smell catalog}
We used ten smells, shown in Table~\ref{tab:freq}. They cover missing context
(vague prompts, unclear references, no stated output format), structure (one
prompt asking for many things, bloat, repetition), contradictions (conflicting
constraints), and tone (excess persona, loaded framing, wrong register).

\subsection{LLM labeling}
Each prompt was sent on its own to \texttt{gpt-oss-120b}~\cite{gptoss} through
the Groq API at temperature 0. The instructions list the ten smells with a short
definition, mark the prompt with fixed start and end markers, and say whether the
prompt is the first turn or a follow-up, so that a short follow-up such as
``make it shorter'' is not marked vague just because it leans on earlier turns.
The model had to return JSON with a smell name and a one-sentence reason for each
smell found. We checked every answer against a schema and asked again, up to three
attempts in all, when it did not fit. Close spellings of a smell name were mapped
to the catalog name. Prompts longer than 12,000 characters were cut, and the model
was told so. Prompts that still failed were logged and retried in a later pass;
the 493 prompts here are those with a valid label.

\section{Results}
\subsection{RQ1: How common are smells?}
Of the 493 prompts, 352 (71.4\%) had at least one smell and 141 (28.6\%) had none.
166 prompts had one smell, 133 had two, and 53 had three or more. Across all
prompts the mean was 1.25 smells, or 1.76 among prompts with any smell.

Table~\ref{tab:freq} shows how often each smell appeared. Missing context led by a
wide margin: it was found in 58.4\% of prompts, more than the next three smells
put together. Format ambiguity (18.7\%) and ambiguous references (16.6\%) came
next. Tone problems were rare. Loaded framing and excess persona each showed up in
about 2\% of prompts, and the model never flagged a formality mismatch. The
labeler was allowed to invent new smell names but did not do so.

\begin{table}[t]
\caption{Prompts with each smell ($n=493$). A prompt can have several.}
\label{tab:freq}
\centering
\begin{tabular}{lrr}
\toprule
Smell & Prompts & \% \\
\midrule
Vague / Missing Context & 288 & 58.4 \\
Format Ambiguity & 92 & 18.7 \\
Ambiguous References & 82 & 16.6 \\
Prompt Bloat / Convoluted Prompt & 45 & 9.1 \\
Overloaded Prompt & 31 & 6.3 \\
Conflicting Constraints & 31 & 6.3 \\
Unnecessary Repetition & 28 & 5.7 \\
Bias / Loaded Framing & 11 & 2.2 \\
Irrelevant / Excessive Persona & 10 & 2.0 \\
Formality / Audience Mismatch & 0 & 0.0 \\
\midrule
Any smell & 352 & 71.4 \\
\bottomrule
\end{tabular}
\end{table}

Smells tended to come in pairs. Of the 82 prompts with an ambiguous reference, 78
were also marked vague, and 62 prompts had both a vague goal and an unclear format.
Of the 31 overloaded prompts, 22 were also bloated.

\subsection{RQ2: Does length matter?}
Table~\ref{tab:len} splits prompts by length. The pattern is U-shaped. Prompts
under 50 characters had a high smell rate (81.3\%), the rate fell to 61.2\% for
prompts of 200--999 characters, and it rose again to 82.6\% for prompts of 1,000
characters or more. The two ends fail in different ways. Among short prompts,
80.2\% were vague and 27.5\% had an ambiguous reference. Among long prompts the
mix was wider: 36.0\% were bloated, 33.7\% vague, 30.2\% had conflicting
constraints and 26.7\% left the format unclear. Long prompts also carried the
most smells each (1.85 on average).

\begin{table}[t]
\caption{Smell rate by prompt length in characters.}
\label{tab:len}
\centering
\begin{tabular}{lrrr}
\toprule
Length & Prompts & With smell (\%) & Mean smells \\
\midrule
$<$ 50 & 91 & 81.3 & 1.16 \\
50--199 & 195 & 68.2 & 1.12 \\
200--999 & 121 & 61.2 & 1.11 \\
$\geq$ 1,000 & 86 & 82.6 & 1.85 \\
\bottomrule
\end{tabular}
\end{table}

\section{Discussion}
The main lesson is simple: most real prompts do not say enough. People seem to
write to ChatGPT the way they would talk to a colleague who already knows the
task, leaving out the goal, the audience or the shape of the answer. The pairing
of ambiguous references with vague prompts fits this view. A user writes ``fix
this'' and assumes the model can see what ``this'' is.

Length is not a fix in itself. Writing more helps up to a point, but very long
prompts bring their own problems. They bury the request, repeat it, or set rules
that pull against each other. Mid-length prompts, a few sentences with a clear
goal, had the lowest smell rate.

For tool builders, this suggests that a light check before sending, such as
asking ``what format do you want?'' or pointing out an ``it'' with nothing to
refer to, could catch many problems. For people who teach prompting, the gains
are likely to come from basics, such as stating the goal and the expected output,
rather than from advanced tricks.

\section{Threats to Validity}
\textbf{Labels.} The labels come from one LLM and were not checked by people.
The model may be too strict or too lenient, and we did not measure agreement with
human raters. The counts should be read as estimates.
\textbf{Sample.} We used one shard, a 1\% sample and English only, and this
paper covers only the first 493 of 1,003 sampled prompts. These come from two
weeks in July 2023 and were all sent to GPT-3.5-Turbo, so we cannot compare
models yet. WildChat language tags are sometimes wrong; we saw a few Chinese
prompts tagged as English. WildChat users chose to share their chats for free
access, so they may not represent all ChatGPT users.
\textbf{Context.} Each prompt was judged alone, with only a first-turn or follow-up
flag. A prompt that looks vague alone may be clear in its chat.

\section{Conclusion}
In 493 real ChatGPT prompts, seven in ten showed at least one prompt smell, and
missing context was the most common by far. Short prompts tend to be vague, and
very long ones tend to be tangled. Next we will finish labeling all 1,003 prompts,
compare models, have people label a subset to check the LLM labels, and test
whether fixing a smell leads to better answers.

\section*{Data and Code Availability}
The labeled data (\texttt{wildchat\_493\_smells.json}), the summary statistics
and the analysis tool are archived on Zenodo at
\url{https://doi.org/10.5281/zenodo.XXXXXXX}. The prompts come from WildChat and
are shared under its ODC-BY license.

\begin{thebibliography}{9}
\bibitem{wildchat} W.~Zhao, X.~Ren, J.~Hessel, C.~Cardie, Y.~Choi, and Y.~Deng,
``WildChat: 1M ChatGPT interaction logs in the wild,'' in \emph{Proc. ICLR}, 2024.
\bibitem{fowler} M.~Fowler, \emph{Refactoring: Improving the Design of Existing
Code}. Addison-Wesley, 1999.
\bibitem{white} J.~White \emph{et al.}, ``A prompt pattern catalog to enhance
prompt engineering with ChatGPT,'' arXiv:2302.11382, 2023.
\bibitem{promptreport} S.~Schulhoff \emph{et al.}, ``The prompt report: A
systematic survey of prompting techniques,'' arXiv:2406.06608, 2024.
\bibitem{johnny} J.~D. Zamfirescu-Pereira, R.~Y. Wong, B.~Hartmann, and Q.~Yang,
``Why Johnny can't prompt: How non-AI experts try (and fail) to design LLM
prompts,'' in \emph{Proc. CHI}, 2023.
\bibitem{judge} L.~Zheng \emph{et al.}, ``Judging LLM-as-a-judge with MT-Bench and
Chatbot Arena,'' in \emph{Proc. NeurIPS Datasets and Benchmarks}, 2023.
\bibitem{gptoss} OpenAI, ``gpt-oss-120b \& gpt-oss-20b model card,''
arXiv:2508.10925, 2025.
\end{thebibliography}

\end{document}
